\section{Analytic continuation of mirror periods}
\label{app:mirror-periods}

\subsection{Toric coordinates and Frobenius series}

The six generators of the ambient nef cone restricted to
$\wideX_9$ form an integral basis. In the divisor basis
$(D_2,D_5,D_6,D_7,D_8,E)$, write $(H_1,\ldots,H_6)$ as the columns
of
\begingroup\setlength{\arraycolsep}{4pt}
\begin{equation}
 N=\begin{pmatrix}
 0&0&1&1&2&2\\
 1&1&2&4&5&5\\
 0&1&2&4&4&5\\
 0&-1&-1&-3&-3&-3\\
 0&-1&-1&-2&-2&-2\\
 0&0&1&2&2&2
 \end{pmatrix},\qquad \det N=-1.
 \label{eq:x9-mirror-nef}
\end{equation}
\endgroup
In particular $H_1=D_5$, $H_2=D_5+D_4$, $H_3=D_1$,
$H_4=D_0$, and $J_*=\sum_aH_a$. The nonnegative orthant in the
dual coordinates is an outer Mori cone sufficient for the period
expansion; its lattice generators need not all be effective on
the hypersurface.

The facet-interior ray $D_f$ misses the hypersurface. Its mirror
monomial coefficient can be removed by an ambient automorphism
\cite[Section~3.1]{Demirtas:2024Mirror}. Removing this ray and
using the regular triangulation induced by the same heights gives
a complete simplicial fan with $24$ maximal cones. Its nef cone
in the six-dimensional divisor space is exactly the cone generated
by \eqref{eq:x9-mirror-nef}; thus removing $D_f$ has not changed
the chosen chamber on the hypersurface.
The GLSM charge matrix, with columns ordered as
$(D_0,D_1,D_2,D_3,D_4,D_5,D_6,D_7,D_8,E)$, is
\begingroup\setlength{\arraycolsep}{4pt}
\begin{equation}
 Q=\begin{pmatrix}
 0&0&-1&0&-1&1&0&1&0&0\\
 0&0&0&1&1&0&0&0&-1&-1\\
 0&1&0&0&0&0&0&2&-2&1\\
 1&0&-1&0&0&0&0&0&0&1\\
 0&0&1&1&0&0&-1&-1&0&0\\
 0&0&0&-1&0&0&1&0&1&-1
 \end{pmatrix}.
 \label{eq:x9-mirror-glsm}
\end{equation}
\endgroup
Its rows are integral relations among the ten ray vectors;
the augmented relation also has central-monomial coefficient
$-Q_{a*}$, where $Q_{a*}=\sum_iQ_{ai}=(0,0,2,1,0,0)_a$.
For the mirror Laurent polynomial
$f=a_* -\sum_i a_i x^{v_i}$, choose algebraic coordinates
$z_a=\prod_i(a_i/a_*)^{Q_{ai}}$.

For $\mathbf n\in\ZZ_{\geq0}^6$, put
$k_i=\sum_a n_aQ_{ai}$ and $k_*=\sum_i k_i$.
The Frobenius series is built from
\begin{equation}
 c(\mathbf n+\boldsymbol\rho)=
 \frac{\Gamma(1+\sum_a(n_a+\rho_a)Q_{a*})}
      {\prod_i\Gamma(1+\sum_a(n_a+\rho_a)Q_{ai})}.
 \label{eq:x9-mirror-gamma}
\end{equation}
Let $W_0=\sum_{\mathbf n}c(\mathbf n)z^{\mathbf n}$,
$C_a=\sum_{\mathbf n}\partial_{\rho_a}c(\mathbf n+\boldsymbol\rho)|_0
z^{\mathbf n}$, and let $\widehat C_{ab}$ use the second
derivative with its universal $\zeta(2)$ multiple of $c(\mathbf n)$
subtracted. Equivalently, all three rational coefficients are
obtained by differentiating
$c(\mathbf n+\boldsymbol\rho)/c(\boldsymbol\rho)$ at zero.
For example, when all $k_i\geq0$ this rational function is
\begin{equation}
 \frac{\prod_{j=1}^{k_*}(j+\sum_a\rho_aQ_{a*})}
      {\prod_i\prod_{j=1}^{k_i}(j+\sum_a\rho_aQ_{ai})}.
 \label{eq:x9-mirror-rational-coefficient}
\end{equation}
For a negative $k_i=-m$, replace that denominator product by
the numerator product
$\prod_{j=0}^{m-1}(\sum_a\rho_aQ_{ai}-j)$.
Such terms may have zero scalar coefficient but nonzero first or
second derivatives; discarding them would lose the conifold
invariants. The mirror map and magnetic generating functions are
\begin{equation}
 \begin{aligned}
 2\pi i t_a&=\log z_a+C_a/W_0,\\
 \sum_{\mathbf n>0}n_a n^0_{\mathbf n}\operatorname{Li}_2(q^{\mathbf n})
 &=\frac12\sum_{b,c}\kappa^{H}_{abc}
       \left(\frac{\widehat C_{bc}}{W_0}
                  -\frac{C_bC_c}{W_0^2}\right),
 \end{aligned}
 \label{eq:x9-mirror-master}
\end{equation}
where $\kappa^H_{abc}=H_aH_bH_c$ is obtained from
\eqref{eq:x9-full-cubic-data} by the basis change
\eqref{eq:x9-mirror-nef}. This is the standard
Frobenius period construction and multicover extraction
\cite[Section~4.4 and Appendix~A]{Demirtas:2024Mirror}.

\subsection{The local exceptional-curve marking}
Relative to the interior point $p$, the five vertices of the
pentagon have coordinates
\[
 (1,0),\ (0,1),\ (-1,0),\ (0,-1),\ (1,1)
\]
in the saturated lattice generated by $v_3-p,v_4-p$.
The face polynomial is therefore
$a_9+a_3\xi+a_4\eta+a_5/\xi+a_6/\eta+a_8\xi\eta$.
With $u=z_1z_5$, $v=z_2$, $w=z_6$, torus rescaling gives
\begin{equation}
 1+vx+y+uw/x+vw/y+xy=0,
 \qquad
 Y^2=(vx^2+x+uw)^2-4vw x^2(1+x).
 \label{eq:x9-local-elliptic-quartic}
\end{equation}
The second equation follows by clearing $xy$ in the first and
taking the discriminant of the resulting quadratic in $y$.
The small-parameter logarithmic correction is the function $H$ below:
$q_6=w e^{-H}$ and $q_2=v e^{-H}$. Its terminating coefficient
identity fixes the local endpoint and the exceptional-curve marking.

\paragraph{Exact endpoint and the two vanishing charges.}
Keep the local mass parameters variable and denote the logarithmic
correction by
\begin{equation}
 H(u,v,w)=\sum_{\substack{0\leq m\leq\ell\leq k\leq\ell+m\\\ell+k>0}}
 \frac{(-1)^{\ell+k-1}(\ell+k-1)!\,u^mv^\ell w^k}
 {m!(\ell+m-k)!(\ell-m)!(k-m)!(k-\ell)!}.
 \label{eq:x9-local-mass-log}
\end{equation}
Thus $q_6=w e^{-H}$ and $q_2=v e^{-H}$.
At $u=1$, summing over $m$ by Vandermonde's identity gives
\begin{equation}
 [v^\ell]H(1,v,w)=
 -\frac{(2\ell-1)!(2\ell)!}{(\ell!)^4}\,w^\ell
 {}_2F_1(-\ell,2\ell;\ell+1;w),\qquad \ell\geq1.
 \label{eq:x9-local-endpoint-polynomial}
\end{equation}
The terminating Chu--Vandermonde identity
\cite[eq.~15.4.24]{DLMF:Hypergeometric} evaluates the last factor
at $w=1$ as $(1-\ell)_\ell/(\ell+1)_\ell=0$.
This coefficient cancellation determines an analytic germ, not
merely a formal substitution: for $|w|\leq1$, the absolute
coefficient at each $v^\ell$ is bounded by
$(\ell+1)^2 125^\ell$, by the multinomial bound on each summand.
Hence $H(1,v,1)=0$ for sufficiently small $|v|$.

For the real continuation to $v=w=1$, the discriminant
on $u=1$ is
\begin{equation}
 \operatorname{disc}_x\mathcal Q
 =-256v^3w^4(w-1)^2\mathcal P(v,w),\qquad
 \mathcal P=16v^2w+(27w^2-18w-1)v+1.
 \label{eq:x9-local-root-wall-discriminant}
\end{equation}
It has no additional zero in $0<v\leq w<1$. Indeed,
the discriminant of $\mathcal P$ as a quadratic in $v$ is
$(w-1)(9w-1)^3$, which is negative for $1/9<w<1$.
For $0<w\leq1/9$, $\mathcal P$ decreases on $0\leq v\leq w$
and $\mathcal P(w,w)\geq1-w-18w^2\geq2/3$.
At $w=1$, $\mathcal P(v,1)=(4v+1)^2$.
Thus continuation along this triangle connects the small-$v$
germ to the diagonal real branch,
without crossing another singularity. It follows that
\begin{equation}
 H(1,v,1)=0\quad(0<v\leq1),\qquad
 \log q_6(1,1,1)=0.
 \label{eq:x9-local-exact-endpoint}
\end{equation}

The fiber itself factors at $u=w=1$:
\begin{equation}
 x(1+x)y^2+(vx^2+x+1)y+vx
 =(xy+1)(vx+xy+y).
 \label{eq:x9-local-two-components}
\end{equation}
The two rational components meet transversely at
$vx^2-x-1=0$, $y=-1/x$ for positive $v$.
The two neck cycles have the same primitive class, up to orientation,
in the compact elliptic curve: either component with small disks
around its two nodes removed supplies a cylinder between them.
Thus the electromagnetic monodromy is conjugate to $T^2$.
This statement forgets the flavor charge carried by the punctures.

The integral marking can be fixed before the collision. As $v\to0$,
\begin{equation}
 \frac{\operatorname{disc}_x\mathcal Q}{-256u^2v^3w^4}
 =(w-1)(uw-1)+O(v),\qquad q_6=w+O(v).
 \label{eq:x9-local-flop-marking}
\end{equation}
These are the exceptional-curve flop loci for $e_2$ and $e_1$,
respectively, since $q^{\iota_*e_2}=q_6$ and
$q^{\iota_*e_1}=u q_6$. Following these two conifold branches
fixes their primitive charges as
$\gamma_i=(0,0,\iota_*e_i,0)$, in the large-volume marking.
The reindexing $m\mapsto k-m$ in
\eqref{eq:x9-local-mass-log} gives the exact Weyl reflection
\begin{equation}
 H(u,v,w)=H(u^{-1},v,uw),\qquad
 (\log u,\log q_6)\longmapsto(-\log u,\log q_6+\log u),
 \label{eq:x9-local-weyl-period}
\end{equation}
on continuously chosen logarithmic branches. It exchanges the
two exceptional-curve periods, with no D0 shift in this marking.
For $u=1+d$, $v=w$ near $1$, the two singularities split as
\begin{equation}
 w_{e_1}=1-\frac{1+\sqrt5}{2}d+O(d^2),\qquad
 w_{e_2}=1+\frac{\sqrt5-1}{2}d+O(d^2).
 \label{eq:x9-local-mass-splitting}
\end{equation}
\subsection{The compact four-node family}
For the finite-bulk calculation, use the monomial basis
\begin{equation}
 m_x=(2,1,1,0),\quad m_y=(2,1,0,1),\quad
 m_r=(1,0,0,0),\quad m_s=(1,1,0,0).
 \label{eq:x9-bulk-monomial-basis}
\end{equation}
Its determinant is one. The last two rows of its inverse are
$n_r=(1,-1,-1,-1)$ and $n_s=(0,1,-1,-1)$, the inward normals
to the two facets meeting in the pentagon. Their cone is smooth,
and the segment joining them has no interior lattice points.
Thus $(\CC^*)^2\times\CC^2$, with regular coordinates $r,s$,
is a valid chart of a crepant toric compactification. In particular,
the loci $r=0$ or $s=0$ must not be discarded when looking for
threefold singularities.

Write $w=z_6$, $v=z_2$, $a=z_3$, $b=z_4$. A coefficient gauge
for the full Laurent polynomial is
\begin{equation}
 a_*=1,\qquad
 (a_0,\ldots,a_9)=
 (b,aw^2z_5^2,1,v,1,z_1z_5w,vw,(wz_5)^{-1},1,1).
 \label{eq:x9-bulk-coefficient-gauge}
\end{equation}
Substitution in \eqref{eq:x9-mirror-glsm} gives exactly
$(z_1,v,a,b,z_5,w)$. Multiplying the Laurent equation by its
nonvanishing torus monomial and extending to this chart gives
\begin{equation}
 \begin{split}
 \mathcal F={}&1+vx+y+\frac{z_1z_5w}{x}+\frac{vw}{y}+xy
       +r\left(1+\frac{xy}{wz_5}\right)-rs\vphantom{\frac11}\\
       &\hspace{25mm}+br^2s+aw^2z_5^2s^2.
 \end{split}
 \label{eq:x9-full-bulk-chart}
\end{equation}
At $z_1=z_5=w=1$, this is \eqref{eq:x9-finite-bulk-factorization}.
Put $f=1+xy$, $g=1+x^{-1}+vy^{-1}+r$ and
$h=r(br-1)+as$. The four proposed singularities satisfy
$f=g=s=h=0$. At either root of $r(br-1)$,
$h_r=2br-1=\pm1$, while
\begin{equation}
 \det\frac{\partial(f,g)}{\partial(x,y)}
       =x^{-1}+vx,\qquad
 (x^{-1}+vx)^2=(1+r)^2+4v.
 \label{eq:x9-bulk-node-transversality}
\end{equation}
Consequently $(f,g,h,s)$ are analytic coordinates at each point
and the equation is $fg+sh=0$. Equivalently, the Hessian
determinant in $(x,y,r,s)$ is $(1+r)^2+4v$. This proves the
ordinary double point assertion to all orders in the bulk
parameters.

The square charts have coordinates
\begin{equation}
 (\xi_1,\eta,\alpha,\sigma)
   =((ry)^{-1},xy,r^{-1},rs),\qquad
 (\xi_2,\eta,\alpha,\sigma)
   =((rx)^{-1},xy,r^{-1},rs).
 \label{eq:x9-bulk-square-charts}
\end{equation}
Their monomial bases, based at $v_2$, are
$(v_6-v_2,v_7-v_2,(-1,0,0,0),(2,1,0,0))$ and the same list
with $v_6$ replaced by $v_5$. Both are unimodular; all transverse
monomial exponents are nonnegative. For positive $b,v$, the
$r=b^{-1}$ roots satisfy
$x_+\sim(bv)^{-1}$ and $x_-\sim-b$ as $b\to0$.
Their respective limits in these charts are
$(\xi_1,\eta,\alpha,\sigma)=(-v^{-1},-1,0,0)$ and
$(\xi_2,\eta,\alpha,\sigma)=(-1,-1,0,0)$, the two square
critical points. The transverse derivative
$\partial_\sigma\mathcal F_{\rm square}$ at
$\alpha=\sigma=0$ is $b$. Thus the undeformed square-face
critical points themselves are smooth points of the full
threefold at $b\neq0$; the actual nodes have moved into the
adjacent facet orbit.

There is a nonempty open subset of this family with no further
singularities. For example, take
\begin{equation}
 (v,a,b)=(1/2,1/100,1/10).
 \label{eq:x9-bulk-rational-sample}
\end{equation}
In the main chart the critical ideal, saturated by $xy$, is
\begin{equation}
 \bigl(s,\ xy+1,\ r(r-10),\ \tfrac12x^2-(1+r)x-1\bigr)
       \subset\QQ[x^{\pm1},y^{\pm1},r,s].
 \label{eq:x9-bulk-critical-ideal}
\end{equation}
It is reduced of length four. Of the twelve facets of the Newton
polytope, eleven have no torus-critical point at these coefficients.
The remaining facet, with indices $(2,3,4,5,6,7,8,9)$, is $s=0$
and has precisely the critical curve $f=g=0$ already treated.
The only nonsimplicial two-faces are the pentagon and two squares;
the latter are smooth in the full threefold by the derivative $b$.
Every remaining lower-dimensional face has affinely independent
support and nonzero coefficients. These exact algebraic checks
exclude other singularities on a crepant toric resolution at
\eqref{eq:x9-bulk-rational-sample}. Properness and the ordinary
double point condition then give a Zariski-open four-node family.

\paragraph{The relation and the deformation rank.}
The closure $W=\overline{\{s=f=0\}}$ is an irreducible compact
Weil divisor of the singular mirror. Near each node it is the
plane $s=f=0$ in $fg+sh=0$, so its local divisor class is a
generator, not a multiple. On $W$, the curve $g=0$ is rational:
\begin{equation}
 y=-x^{-1},\qquad r=-1-x^{-1}+vx,\qquad
 x^2h=(vx^2-x-1)\bigl(b(vx^2-x-1)-x\bigr).
 \label{eq:x9-bulk-four-intersections}
\end{equation}
The last polynomial has degree four, constant term $b$ and
leading coefficient $bv^2$. Its four simple zeros are all the
nodes. The two missing endpoints of the rational curve are
covered by \eqref{eq:x9-bulk-square-charts}, where the normal
derivative is nonzero. Thus $W$ meets no additional singularity
for generic bulk parameters. Cutting out four local balls and
transporting the complement to a nearby smooth mirror produces
a four-chain whose boundary is the sum of the four oriented
vanishing spheres. This is the usual conifold four-chain
construction \cite[Section~3]{Greene:1995hu}; it gives
a relation with four unit coefficients in absolute value.
Blowing up $W$ gives a projective small resolution of all four
nodes and is an isomorphism elsewhere.

At a node, parameter differentiation gives
\eqref{eq:x9-finite-bulk-smoothing}; derivatives in $v,a,b$ vanish.
For the two roots at $r=0$, write their inverse $x$-coordinates
as $\rho\neq\rho'$. Either root at $r=b^{-1}$, with inverse
coordinate $\tau$, gives the minor
\begin{equation}
 \det\begin{pmatrix}
 \rho&\rho&-1\\ \rho'&\rho'&-1\\
 \tau&\tau+b^{-1}&-1
 \end{pmatrix}=\frac{\rho-\rho'}{b}\neq0.
 \label{eq:x9-bulk-rank-three}
\end{equation}
The leading period of a vanishing sphere is a nonzero multiple
of its smoothing parameter. Thus these derivatives imply at least
three independent homology classes. The four-chain relation gives
at most three, proving rank three and uniqueness of the primitive
relation.

One can also check the relative coefficients from residues.
Up to a common sign, the holomorphic form in this chart is
$\operatorname{Res}(dx\,dy\,dr\,ds/(xy\mathcal F))$.
Set $D=x^{-1}+vx$ at a node, so
$x^{-1}=(D-1-r)/2$ and $D=\pm\sqrt{(1+r)^2+4v}$.
After removing a common period normalization, its smoothing
covector is
\begin{equation}
 \frac{(x^{-1},x^{-1}+r,-1)}{D(2br-1)}.
 \label{eq:x9-bulk-residue-covectors}
\end{equation}
At fixed $r$, the sum over the two signs of $D$ is
$(1,1,0)/(2br-1)$. The sums at $r=0$ and $r=b^{-1}$ therefore
cancel, in agreement with the integral relation. This residue
calculation is a first-order consistency check; the four-chain
argument, not equality of four derivatives, establishes the
homology relation. To identify the embedding in the large-volume
charge lattice, one must also fix the electric periods and their
normalization. We do this next on a marked small-bulk neighborhood.

\subsection{Electric continuation and exact flat coordinates}
Denote \eqref{eq:x9-finite-bulk-locus} by $\Sigma$ and keep
$(v,a,b)=(z_2,z_3,z_4)$ sufficiently small and nonzero. At fixed
spectator degrees $(\ell,m,n)$, write the summation degree in
\eqref{eq:x9-mirror-gamma} as $(i,\ell,m,n,j,k)$. The ten
denominator arguments, with the constant $1$ omitted, are
\begin{equation}
 \begin{split}
 (k_0,\ldots,k_9)=(&n,m,j-i-n,\ell+j-k,\ell-i,i,\\
                  &k-j,i+2m-j,k-\ell-2m,m+n-\ell-k),
 \end{split}
 \label{eq:x9-bulk-gamma-arguments}
\end{equation}
and $k_*=2m+n$. Two negative arguments make the scalar and every
first derivative vanish. Requiring all but at most one argument
to be nonnegative gives
\begin{equation}
 \begin{aligned}
 0\leq i&\leq\max(\ell,m-\ell),\\
 0\leq j&\leq\max(\ell+2m,m+n-\ell),\\
 0\leq k&\leq\max(2\ell+2m,m+n-\ell).
 \end{aligned}
 \label{eq:x9-bulk-first-jet-bounds}
\end{equation}
These bounds follow by choosing the possible negative argument
in \eqref{eq:x9-bulk-gamma-arguments} and solving the remaining
linear inequalities. For example, if $i>\ell$, then $\ell-i$
is the only negative argument, so $i+n\leq j\leq k\leq m+n-\ell$
and hence $i\leq m-\ell$.
Thus each coefficient in $(v,a,b)$ contains only finitely many
local degrees, all of which can be summed exactly.

This regrouping also proves convergence beyond $|z_1|,|z_5|,|z_6|<1$.
Put $d=\ell+m+n$. The bounds imply $i+j+k\leq5d$,
$|k_i|\leq3d$ and $k_*\leq2d$. With one negative argument
$-s$, the positive arguments sum to $k_*+s\leq5d$;
the nonzero derivative coefficient is a bounded charge-matrix
entry times $k_*!(s-1)!/\prod_{k_i\geq0}k_i!$.
The multinomial bound therefore bounds this coefficient by a
constant times $10^{5d}$. With no negative argument the same
bound holds, up to polynomial factors from harmonic numbers.
There are only polynomially many local summands at each $d$.
For any fixed $M>1$, all seven series $W_0,C_1,\ldots,C_6$
consequently converge normally on
$|z_1|,|z_5|,|z_6|\leq M$ when $|v|,|a|,|b|$ are sufficiently
small. At $(v,a,b)=0$, they are $W_0=1$ and $C_a=0$.
In a smaller neighborhood $W_0$ is nonzero.

The overlap with the original large-volume series fixes their
continuation. Choose logarithms by continuation from positive
small $z_a$, without coordinate winding. Then all seven electric
periods
\begin{equation}
 \Pi_0=W_0,\qquad
 \Pi_a=\frac{W_0\log z_a+C_a}{2\pi i}=W_0t_a
 \quad(1\leq a\leq6)
 \label{eq:x9-bulk-electric-periods}
\end{equation}
are holomorphic and single-valued near $\Sigma$ at nonzero small
spectator parameters. This is analytic continuation of the marked
periods, not substitution into a finite total-degree expansion.

\paragraph{The exact flat-coordinate equations.}
To evaluate the three local first derivatives on $\Sigma$, put
\begin{equation}
 B_N(t)=\frac1{\Gamma(1+t)\Gamma(1+N-t)}.
 \label{eq:x9-bulk-binomial-kernel}
\end{equation}
Apart from $(2m+n)!/(m!n!)$, the fixed-spectator coefficient
factors into four such kernels. Their orders, integer indices
and shifts under $(\rho_1,\rho_5,\rho_6)$ are
\begin{equation}
 \begin{array}{c|c|c}
 N_\alpha&I_\alpha&\epsilon_\alpha\\\hline
 \ell&i&\rho_1\\
 2m-n&j-i-n&\rho_5-\rho_1\\
 \ell&k-j&\rho_6-\rho_5\\
 n-m-2\ell&m+n-\ell-k&-\rho_6
 \end{array},\qquad
 \begin{gathered}
 \sum_\alpha N_\alpha=m,\\
 \sum_\alpha I_\alpha=m-\ell,\\
 \sum_\alpha\epsilon_\alpha=0.
 \end{gathered}
 \label{eq:x9-bulk-four-binomial-pairs}
\end{equation}
The first-order sum may be extended to all integer $i,j,k$:
outside the nonnegative orthant there are at least two negative
arguments. Indeed, $i<0$ would require both $j\geq2m$ and
$j\leq i+2m$ if every other argument were nonnegative.
With $i\geq0$, $j<0$ similarly requires $j\geq2m$; with
$i,j\geq0$, $k<0$ makes both $k-j$ and $k-\ell-2m$ negative.
The bounds above show that the first-order convolution has finite
support, so no regularization of an infinite sum is involved.

For $N\geq0$, the sequence $B_N(k)$ at integer $k$ is
$\binom Nk/N!$, supported on $0\leq k\leq N$. The Gamma
recurrence gives
\begin{equation}
 B_N(t)=\frac{B_{N-1}(t)+B_{N-1}(t-1)}{N},\qquad
 B_N'=B_N*B_0',\qquad
 B_0'(k)=\begin{cases}(-1)^k/k,&k\neq0,\\0,&k=0,\end{cases}
 \label{eq:x9-bulk-binomial-convolution}
\end{equation}
where the recurrence is for $N\geq1$, and the convolution identity
is an identity of integer-indexed sequences. If all four orders
are nonnegative, differentiating any one factor in their
convolution gives the same result. The zero sum of the shifts in
\eqref{eq:x9-bulk-four-binomial-pairs} therefore kills every local
first derivative.
If one order is $-s<0$, then $B_{-s}(k)=0$ for all integers and
\begin{equation}
 B_{-s}'(k)=(-1)^{s+k}\prod_{r=1}^{s-1}(k+r).
 \label{eq:x9-bulk-negative-binomial}
\end{equation}
The other three kernels combine into a binomial kernel of degree
$m+s\geq s$, which annihilates this alternating polynomial of
degree $s-1$ by finite differences. If two orders are negative,
the first derivative is already zero. These cases prove, to all
orders in the spectator parameters,
\begin{equation}
 C_1|_\Sigma=C_5|_\Sigma=C_6|_\Sigma=0,\qquad
 t_1|_\Sigma=t_5|_\Sigma=t_6|_\Sigma=0.
 \label{eq:x9-bulk-exact-flat-locus}
\end{equation}
Normal convergence makes these analytic identities. Ordinary
Vandermonde convolution also gives
\begin{equation}
 [v^\ell a^m b^n]W_0|_\Sigma
 =\frac{(2m+n)!}
 {n!(\ell!)^3(m-\ell)!(2m-n)!(n-m-2\ell)!},
 \label{eq:x9-bulk-restricted-fundamental}
\end{equation}
with a reciprocal negative factorial interpreted as zero.
Finally, on $\Sigma$,
\begin{equation}
 2\pi i(t_2,t_3,t_4)=(\log v,\log a,\log b)+O(v,a,b).
 \label{eq:x9-bulk-spectator-independence}
\end{equation}
Their Jacobian with respect to the spectator
logarithms tends to the identity, so these three flat coordinates
are independent for sufficiently small nonzero bulk parameters.

\paragraph{The integral charges and their unit normalization.}
Let $\delta$ be any of the four vanishing spheres and let
$(p^0,P,q,q_0)$ be its charge in the continued large-volume frame.
Picard--Lefschetz monodromy around its individual discriminant
branch acts by
\begin{equation}
 \Pi_\eta\longmapsto\Pi_\eta+
       \langle\eta,\delta\rangle\Pi_\delta.
 \label{eq:x9-bulk-electric-invariance}
\end{equation}
The period $\Pi_\delta$ is not the zero germ: its first derivative
in the smoothing direction is nonzero. All seven electric periods
\eqref{eq:x9-bulk-electric-periods} are single-valued, so each
pairs trivially with $\delta$. They span the electric Lagrangian
sublattice in the large-volume marking; hence $p^0=P=0$.
Dividing the vanishing period by $W_0$ and restricting to $\Sigma$
then gives
$q_2t_2+q_3t_3+q_4t_4-q_0=0$.
Independence of the spectator flat coordinates forces
$q_2=q_3=q_4=q_0=0$. Thus every sphere lies in
$\ZZ\beta_1\oplus\ZZ\beta_5\oplus\ZZ\beta_6$, with no magnetic
or D0 component.

To fix its integer multiplicity as well as its direction, normalize
the toric residue by $(2\pi i)^{-3}$, so that $W_0\to1$ at
large volume. In the node model $uv+st=\mu$, the vanishing
$S^3$ is a two-torus fibration over $q=uv\in[0,\mu]$.
Integrating the residue gives $\pm(2\pi i)^2\mu$ before this
normalization, and hence $\pm\mu/(2\pi i)$ afterwards. The
covectors \eqref{eq:x9-bulk-residue-covectors} therefore have the
same unit normalization as the limiting flat differentials
$dt_a=d\log z_a/(2\pi i)$.
For positive $(v,a,b)\to0$, write $x_+>0>x_-$ at each of the
two values of $r$. With the indicated choice of orientations,
the limiting covectors and the resulting exact charges are
\begin{equation}
 \begin{array}{c|c|c|c}
 r&x& (q_1,q_5,q_6)&q\\\hline
 b^{-1}&x_+&(0,1,0)&C_1=\beta_5\\
 b^{-1}&x_-&(1,0,0)&C_2=\beta_1\\
 0&x_-&(1,1,1)&\iota_*e_1=\beta_1+\beta_5+\beta_6\\
 0&x_+&(0,0,1)&\iota_*e_2=\beta_6
 \end{array}.
 \label{eq:x9-bulk-unit-charge-marking}
\end{equation}
The third row reverses the common residue orientation. The
integer charges are locally constant, so their limiting values
identify them exactly at nonzero small bulk parameters, not only
at the boundary. Their span is saturated and their unique
primitive relation is $\iota_*e_1-\iota_*e_2=C_1+C_2$.

There is no homological torsion ambiguity in this identification.
The vertices of the original fan polytope contain the standard
basis of $\ZZ^4$. The polar polytope has the four vertices
\begin{equation}
 (-1,2,1,-1),\quad(1,-1,-1,-1),\quad
 (0,-1,-1,1),\quad(-1,-1,-1,4),
 \label{eq:x9-bulk-polar-lattice}
\end{equation}
whose determinant has absolute value one. In both polytopes the
vertices alone therefore generate the full lattice. The toric
fundamental-group and Brauer-group formulas
\cite[Corollaries~1.9 and~3.9]{Batyrev:2005Integral}
give trivial fundamental and Brauer groups for the resolved
threefold and its smooth mirror; in particular the mirror
$H_3$ has no torsion.

The marking is fixed by the convergent continuation just described,
on a nonempty small-bulk part of the four-node family. This proof
does not compute transport to the particular point
\eqref{eq:x9-bulk-rational-sample} or around arbitrary global
loops. Further continuation must transport the charge frame as
well as the periods. Nor does electric monodromy determine the
regular parts of the magnetic periods or bound-state stability.

\subsection{Regular magnetic periods and their Abel limit}
We now determine the restricted regular terms. For $P=\sum_a p^aH_a$,
put $\kappa_{Pbc}=PH_bH_c$, $g_b=C_b/W_0$ and
\begin{equation}
 D_P=\frac12\sum_{b,c}\kappa_{Pbc}\widehat C_{bc},\qquad
 \Psi_P=\frac1{4\pi^2}
       \left(\frac{D_P}{W_0}
                    -\frac12\sum_{b,c}\kappa_{Pbc}g_bg_c\right).
 \label{eq:x9-endpoint-magnetic-master}
\end{equation}
This is the magnetic normalization of \eqref{eq:x9-mirror-master}. With at least three
negative denominator arguments, every second derivative vanishes.
With two negative arguments $k_i,k_j$, its contraction is
\begin{equation}
 (PD_iD_j)\,
 \frac{(-1)^{-k_i-k_j-2}(2m+n)!(-k_i-1)!(-k_j-1)!}
      {\prod_{r\neq i,j} k_r!}.
 \label{eq:x9-two-zero-magnetic-coefficient}
\end{equation}
All other arguments in this formula are nonnegative. The possible
unbounded local-degree rays, at fixed $(\ell,m,n)$, have negative
pairs $(2,4),(6,7),(4,9),(3,9),(4,6),(7,9)$ in the zero-based
column ordering of \eqref{eq:x9-mirror-glsm}. The last two have
$D_4D_6=D_7E=0$ on the hypersurface and contribute nothing.
The first four have $D_iD_j=C_2,C_1,\iota_*e_1,\iota_*e_2$,
respectively. Every other nonzero sector has $i,j,k\leq4d$,
where $d=\ell+m+n$; this follows from the remaining linear
inequalities in \eqref{eq:x9-bulk-gamma-arguments}.

The four unbounded sectors are one-dimensional after fixing their
transverse indices. The following table gives all such indices for
the unbounded integer $h$.
An empty index interval contributes zero.
\begingroup\setlength{\arraycolsep}{3pt}
\begin{equation}
 \begin{array}{c|c|c}
 D_iD_j&(i,j,k)&\text{transverse indices}\\\hline
 C_2&(h,j,k)&\ell+2m\leq k\leq m+n-\ell,
                  \ k-\ell\leq j\leq k\\
 C_1&(i,h,k)&0\leq i\leq\ell,
                  \ \ell+2m\leq k\leq m+n-\ell\\
 \iota_*e_1&(h,h+p,h+p+q)&n\leq p\leq2m,
                  \ 0\leq q\leq\ell\\
 \iota_*e_2&(i,j,h)&0\leq i\leq\ell,
                  \ i+n\leq j\leq i+2m
 \end{array}.
 \label{eq:x9-magnetic-tail-lines}
\end{equation}
\endgroup
Their respective lower bounds are
\begin{equation}
 \begin{aligned}
 h_{C_2}&=\max(\ell+1,j-n+1,j-2m),\\
 h_{C_1}&=\max(k+1,i+2m+1,i+n),\\
 h_{e_1}&=\max(\ell+1,m+n-\ell-p-q+1,\ell+2m-p-q),\\
 h_{e_2}&=\max(\ell+j+1,m+n-\ell+1,\ell+2m).
 \end{aligned}
 \label{eq:x9-magnetic-tail-lower-bounds}
\end{equation}
Take a common cutoff $h_0=8d+2$, exceeding all these lower bounds.
Below it the coefficient sum is finite. Above it, sum the
transverse indices before summing $h$. To display the resulting
rational functions, define
\begin{equation}
 N_-=n-m-2\ell,\quad N_+=2m-n,\quad
 K=\frac{(2m+n)!}{m!n!},\quad
 x^{\underline r}=\prod_{j=0}^{r-1}(x-j).
 \label{eq:x9-magnetic-tail-notation}
\end{equation}
The two nodal tails are zero when $N_-<0$, and the two
exceptional-curve tails are zero when $N_+<0$. Otherwise write
the scalar tail for charge $\gamma$ as $A_\gamma R_\gamma(h)$,
where
\begin{equation}
 \begin{aligned}
 A_{C_2}&=\frac{K(-1)^{n+\ell}
             (n-2m-1)^{\underline{\ell+N_-}}}{\ell!N_-!},\\
 A_{C_1}&=\frac{K(-1)^\ell(-\ell-1)^{\underline{N_-}}
             (n-2m-1)^{\underline\ell}}{\ell!N_-!},\\
 A_{e_1}&=\frac{K(-1)^{m+\ell+N_+}
             (m+2\ell-n-1)^{\underline{\ell+N_+}}}{\ell!N_+!},\\
 A_{e_2}&=\frac{K(-1)^m
             (-\ell-1)^{\underline{\ell+N_+}}}{\ell!N_+!},
 \end{aligned}
 \label{eq:x9-magnetic-tail-amplitudes}
\end{equation}
and
\begin{equation}
 \begin{aligned}
 R_{C_2}(h)=R_{e_1}(h)
  &=\frac{\Gamma(h-\ell)\Gamma(h+\ell-m)}{\Gamma(h+1)^2},\\
 R_{C_1}(h)
  &=\frac{\Gamma(h-m-n+\ell)\Gamma(h-\ell-2m)}
          {\Gamma(h-2m+1)\Gamma(h-n+1)},\\
 R_{e_2}(h)
  &=\frac{\Gamma(h-2m-2\ell)\Gamma(h-m-n+\ell)}
          {\Gamma(h-n+1)\Gamma(h-\ell-2m+1)}.
 \end{aligned}
 \label{eq:x9-magnetic-rational-tails}
\end{equation}
These formulas follow by the finite-difference identity
\begin{equation}
 \sum_{j=0}^r(-1)^j\binom rj
       \frac{\Gamma(x-j+\lambda)}{\Gamma(x-j+1)}
 = (\lambda-1)^{\underline r}
       \frac{\Gamma(x-r+\lambda)}{\Gamma(x+1)},
 \label{eq:x9-magnetic-finite-difference}
\end{equation}
and its forward version, applied to the factorial pairs in
\eqref{eq:x9-magnetic-tail-lines}. The identity itself follows
inductively from the Gamma recurrence.

Every nonzero rational tail is $O(h^{-m-2})$. Its poles are
integers below $h_0$, of order at most two. Partial fractions
therefore evaluate its sum exactly in $\QQ+\QQ\zeta(2)$:
the simple-pole residues sum to zero, giving rational harmonic
number differences, while double poles give
$\zeta(2)-\sum_{j=1}^{h_0-r-1}j^{-2}$ for a pole at $r$.
Together with the finite terms, this determines every
$[v^\ell a^m b^n]D_P|_\Sigma$ without truncating local degrees.

\paragraph{Convergence and the endpoint limit.}
There are constants $c_0,c_1$ independent of $(\ell,m,n)$ such
that the absolute value of each fully summed coefficient is at
most $c_0c_1^d$. For the finite part, all Gamma arguments are
$O(d)$, and the factorial multinomial bound used above applies
also to \eqref{eq:x9-two-zero-magnetic-coefficient}; harmonic
numbers and the number of terms add only polynomial factors.
For the tail, all shifts in \eqref{eq:x9-magnetic-rational-tails}
have absolute value at most $4d+1$. Comparing each linear factor
at $h\geq h_0$ with its value at $h_0$ gives, for a constant $c$,
\begin{equation}
 |A_\gamma R_\gamma(h)|
 \leq c^d|A_\gamma R_\gamma(h_0)|
                  (h_0/h)^{m+2}.
 \label{eq:x9-magnetic-tail-bound}
\end{equation}
The value at $h_0$ is bounded by the finite transverse factorial
sum. Summing the last bound costs only $O(h_0)$ and another
exponential factor in $d$. This proves normal convergence of
the restricted bulk series.

To check that the sum is the continued period value, rather than
an unjustified specialization, first put $z_1=z_5=z_6=\rho<1$.
Each line in \eqref{eq:x9-magnetic-tail-lines} carries its full
weight $\rho^{i+j+k}$, and its $h$ sum converges. In a transverse
finite difference, replacing a difference by the extra weight
costs a factor $O(1-\rho)$ and loses one inverse power of $h$.
Consequently the difference from the endpoint tail is a finite
sum of terms bounded, at fixed spectator degree, by
\begin{equation}
 (1-\rho)^s\sum_{h\geq h_0}
      \rho^{ch}h^{-m-2+s},\qquad s\geq1,\quad c\in\{1,3\}.
 \label{eq:x9-magnetic-abel-bound}
\end{equation}
Here is a uniform version of this estimate. If $E_h$ is the unit
shift in a transverse difference, the weighted difference satisfies
\begin{equation}
 (1-uE_h)^r=\sum_{s=0}^r\binom rs(1-u)^s
                     u^{r-s}(1-E_h)^{r-s}.
 \label{eq:x9-weighted-difference}
\end{equation}
The transverse orders in the four sectors have total at most $3d$;
their weights $u$ are fixed positive powers of $\rho$. For
$1/2\leq\rho<1$, all bounded shifts of the residual weight cost
only an exponential factor in $d$. Applying the Gamma recurrence
to each remaining difference, and comparing its linear factors at
$h_0=8d+2$, bounds its contribution by
\begin{equation}
 C^d\epsilon^s\sum_{h\geq h_0}\rho^{ch}
       (h/h_0)^p,\qquad
 \epsilon=1-\rho,\quad p=s-m-2,\quad 0\leq s\leq3d.
 \label{eq:x9-uniform-weighted-tail}
\end{equation}
Polynomial factors in $d$ are absorbed in $C^d$ (the degree-zero
case is immediate). The coefficient at $h_0$ is bounded by the
same finite factorial sums used in \eqref{eq:x9-magnetic-tail-bound};
the binomial factors in \eqref{eq:x9-weighted-difference} add at
most $2^{3d}$.
For $p\geq0$, comparison with an exponential moment bounds
\eqref{eq:x9-uniform-weighted-tail}, after increasing $C$, by
\begin{equation}
 C^d\frac{p!}{h_0^p}\epsilon^{s-p-1}
 =C^d\frac{p!}{h_0^p}\epsilon^{m+1}.
\end{equation}
Since $p\leq3d<h_0$, $p!/h_0^p\leq1$. For $p=-1$ the
bound is $C^dh_0\epsilon^s(1+|\log\epsilon|)$; for $p\leq-2$
it is $C^dh_0\epsilon^s$. Every term with $s\geq1$ tends to
zero and has an exponential-in-$d$ bound uniform in $\rho$.
The $s=0$ term has the summable endpoint decay $h^{-m-2}$ and
converges by dominated convergence. The number of transverse
terms is polynomial in $d$, so the same majorant applies to their
sum. On a smaller spectator polydisc it is summable also in
$(\ell,m,n)$, justifying interchange of the Abel limit and the
spectator series. Its overlap with the large-volume Frobenius series fixes
the continuation marking. The magnetic periods have the usual
transverse vanishing-period logarithms, but their restrictions
to $\Sigma$ are the analytic functions just obtained.

For completeness, the constant and linear coefficients of
$(D_{H_1},\ldots,D_{H_6})|_\Sigma$ are
\begin{equation}
 \begin{aligned}
 [1]D_H&=2\zeta(2)(1,0,0,0,1,1),\\
 [v]D_H=[a]D_H&=(1,0,0,0,1,2),\\
 [b]D_H&=(0,0,0,10,10,10).
 \end{aligned}
 \label{eq:x9-magnetic-endpoint-coefficients}
\end{equation}
All these coefficient sums are real for positive spectator parameters.

\paragraph{The four conifold constants.}
An independent boundary calculation checks the coefficient and sign.
At spectator degree zero, all first Gamma derivatives vanish. The
nonzero local contributions to the cubic Taylor coefficient
$\kappa_{abc}\partial_a\partial_b\partial_c/6$ lie on the four
rays $hC_1$, $hC_2$, $h\iota_*e_1$ and $h\iota_*e_2$, with $h\geq1$.
The other two possible two-zero sectors have zero divisor
intersection; every three-zero sector has zero triple intersection.
On each active ray the two reciprocal Gamma zeros supply $h^{-2}$.
Writing $\mathsf H_h=\sum_{j=1}^h j^{-1}$, the remaining
logarithmic derivative contracts to
$-2(\mathsf H_h-\mathsf H_{h-1})=-2/h$. Thus the four complete
ray sums give
\begin{equation}
 4\sum_{h\geq1}\frac{-2}{h^3}=-8\zeta(3)
 \label{eq:x9-d6-four-tails}
\end{equation}
in the cubic Gamma coefficient, or
$-4\zeta(3)/(2\pi i)^3$ in the instanton prepotential. These are
the four $\operatorname{Li}_3(1)$ contributions with unit genus-zero
index. Their sum compensates the change of the perturbative Euler
term exactly. The boundary calculation checks the constant; the
polarization and reality argument proves the full period identity.
